%==============================================================================
%  CAHIER DES CHARGES — ASSURANCE NATION
%  ENSPY — 3GI Groupe 9 — Conception des Systemes d'Information (CSI)
%  Compilation : pdflatex (x2) Cahier_des_Charges_ASSURANCE_NATION.tex
%  Dependances : tikz, pgfgantt, longtable, booktabs, hyperref, fancyhdr...
%==============================================================================
\documentclass[11pt,a4paper]{article}

\usepackage[utf8]{inputenc}
\usepackage[T1]{fontenc}
\usepackage[french]{babel}
\usepackage{lmodern}
\usepackage[a4paper,margin=2.5cm]{geometry}
\usepackage{graphicx}
\usepackage[table]{xcolor}
\usepackage{array}
\usepackage{longtable}
\usepackage{booktabs}
\usepackage{tabularx}
\usepackage{multirow}
\usepackage{enumitem}
\usepackage{titlesec}
\usepackage{fancyhdr}
\usepackage{tikz}
\usepackage{pgfgantt}
\usepackage{caption}
\usepackage{float}
\usetikzlibrary{shapes.geometric,arrows.meta,positioning,fit,backgrounds}
\usepackage[hidelinks,colorlinks=true,linkcolor=anbleu,urlcolor=anbleu,citecolor=anbleu]{hyperref}

%------------------------------------------------------------------ Couleurs
\definecolor{anbleu}{HTML}{2E3192}   % bleu institutionnel
\definecolor{anjaune}{HTML}{FFE100}  % jaune
\definecolor{angris}{HTML}{555555}
\definecolor{anligne}{HTML}{2E3192}

%------------------------------------------------------------------ Titres
\titleformat{\section}
  {\normalfont\Large\bfseries\color{anbleu}}{\thesection}{1em}{}
\titleformat{\subsection}
  {\normalfont\large\bfseries\color{anbleu}}{\thesubsection}{1em}{}
\titleformat{\subsubsection}
  {\normalfont\normalsize\bfseries\color{angris}}{\thesubsubsection}{1em}{}

%------------------------------------------------------------------ En-tetes
\pagestyle{fancy}
\fancyhf{}
\setlength{\headheight}{14pt}
\fancyhead[L]{\small\itshape Cahier des Charges — ASSURANCE NATION}
\fancyhead[R]{\small\itshape ENSPY · 3GI Groupe 9}
\fancyfoot[C]{\thepage}
\renewcommand{\headrulewidth}{0.4pt}
\renewcommand{\footrulewidth}{0pt}

%------------------------------------------------------------------ Tableaux
\newcolumntype{L}[1]{>{\raggedright\arraybackslash}p{#1}}
\renewcommand{\arraystretch}{1.25}

%------------------------------------------------------------------ Exigence (boite ISO 29148)
\newcommand{\exig}[7]{%
\par\noindent
\begin{tabularx}{\linewidth}{|>{\bfseries}p{3.1cm}|X|}
\hline
\rowcolor{anbleu!10}\multicolumn{2}{|l|}{\textbf{\color{anbleu}#1} — #2}\\\hline
Description & #3 \\\hline
Priorité (MoSCoW) & #4 \\\hline
Complexité / Effort & #5 \\\hline
Critères d'acceptation & #6 \\\hline
Méthode de vérification & #7 \\\hline
\end{tabularx}\par\vspace{0.6em}}

\begin{document}

%==============================================================================
%  PAGE DE COUVERTURE  (sans mention de version)
%==============================================================================
\begin{titlepage}
\thispagestyle{empty}
\begin{tikzpicture}[remember picture,overlay]
  % --- bandes decoratives ---
  \fill[anbleu]  (current page.north west) rectangle ++(7.0,-3.6);
  \fill[anjaune] (current page.north west) ++(0,-3.6) rectangle ++(3.4,-0.5);
  \fill[anbleu]  (current page.south east) rectangle ++(-7.4,3.6);
  \fill[anjaune] (current page.south west) ++(0,2.0) rectangle ++(2.6,1.2);

  % --- en-tete (dans la bande bleue) ---
  \node[anchor=north west,white,font=\bfseries\large,align=left,text width=6.4cm]
    at ([shift={(0.5,-0.6)}]current page.north west)
    {CONCEPTION DES\\SYSTÈMES\\D'INFORMATION};
  \node[anchor=north west,anbleu,font=\bfseries\Large]
    at ([shift={(0.5,-4.6)}]current page.north west) {3GI --- Groupe 9};

  % --- titre central ---
  \node[anchor=center,align=center] at ([yshift=2.2cm]current page.center) {%
    {\color{anbleu}\fontsize{40}{44}\selectfont\bfseries CAHIER DES CHARGES}\\[0.5cm]
    {\fontsize{22}{26}\selectfont\bfseries Plateforme \textsc{Assurance Nation}}\\[0.35cm]
    {\large\itshape Gestion des consultations médicales et des remboursements}\\[0.15cm]
    {\large\itshape d'un organisme de sécurité sociale}};

  % --- etablissement + supervision (bas gauche, sur fond blanc) ---
  \node[anchor=south west,align=left] at ([shift={(0.5,4.4)}]current page.south west) {%
    {\color{anbleu}\bfseries\large ÉTABLISSEMENT}\\[0.15cm]
    \normalsize École Nationale Supérieure Polytechnique de Yaoundé (ENSPY)\\[0.7cm]
    {\color{anbleu}\bfseries\large SOUS LA SUPERVISION DE :}\\[0.15cm]
    \normalsize Dr. Anne Marie CHANA\\
    \normalsize Dr. Jaures Styve KAMENI};

  % --- annee academique (bas droit, en blanc sur la bande bleue) ---
  \node[anchor=south east,white,font=\bfseries\large,align=right]
    at ([shift={(-0.6,0.8)}]current page.south east) {Année Académique :\\2025--2026};
\end{tikzpicture}
\end{titlepage}

%==============================================================================
%  TABLE DES MATIERES
%==============================================================================
\tableofcontents
\thispagestyle{empty}
\clearpage

%==============================================================================
%  PREAMBULE
%==============================================================================
\section*{Préambule}
\addcontentsline{toc}{section}{Préambule}
\markboth{Préambule}{Préambule}

Le présent document constitue le \textbf{cahier des charges} de la plateforme
logicielle \textsc{Assurance Nation}. Il est rédigé conformément aux principes
structurels des normes \textbf{IEEE~830} et \textbf{ISO/IEC/IEEE~29148} relatives
à l'ingénierie des exigences. Son objet est de formaliser, sans ambiguïté et de
manière vérifiable, l'ensemble des besoins métiers, des exigences fonctionnelles
et non-fonctionnelles, ainsi que les contraintes d'architecture, de sécurité,
d'hébergement, de méthodologie et de cadre juridique qui régissent la conception,
la réalisation et la livraison du système.

Ce cahier des charges fait suite au \emph{Cahier d'Analyse} du projet tutoré
(modélisation UML : diagrammes de contexte, de packages, de classes, de cas
d'utilisation, d'activités et de séquences) et en consolide les conclusions sous
une forme contractuelle et normée.

%==============================================================================
%  SECTION 1
%==============================================================================
\setcounter{section}{0}
\section{Introduction et Contexte Stratégique}

La présente section a pour fonction de \textbf{contextualiser le projet}. Les
développeurs et architectes logiciels ne peuvent concevoir une solution optimale
s'ils ignorent la finalité métier du système. Elle aligne donc la compréhension de
l'ensemble des acteurs (maîtrise d'ouvrage, maîtrise d'œuvre, encadrement
pédagogique) autour d'une vision commune.

\subsection{Définition du Problème et État de l'Art}

\subsubsection{Description de l'existant}
\textsc{Assurance Nation} est destiné à un \textbf{organisme de sécurité sociale}
qui souhaite mettre en place une application centralisée pour la gestion des
patients (assurés), de leurs médecins traitants, des médecins spécialistes, ainsi
que des remboursements médicaux liés aux consultations.

À l'état initial, l'organisme repose sur des \textbf{processus essentiellement
manuels et hétérogènes} :
\begin{itemize}[leftmargin=1.4em]
  \item l'inscription des assurés et des médecins s'effectue sur supports papier
  ou tableurs non centralisés ;
  \item les feuilles de maladie sont remplies à la main, puis transmises
  physiquement entre le médecin, l'assuré et l'agent de sécurité sociale ;
  \item le calcul des remboursements (100~\% pour une consultation chez un
  généraliste, 80~\% chez un spécialiste) est réalisé manuellement, ce qui est
  source d'erreurs ;
  \item le suivi du parcours de soins (consultations, prescriptions,
  orientations vers spécialistes) n'est pas tracé de bout en bout ;
  \item aucune piste d'audit fiable ne permet de savoir qui a modifié quelle
  donnée et à quel moment.
\end{itemize}

\subsubsection{Points de friction et limites}
\begin{itemize}[leftmargin=1.4em]
  \item \textbf{Perte de productivité} : ressaisies multiples, recherches
  documentaires fastidieuses, délais de traitement des dossiers de remboursement.
  \item \textbf{Taux d'erreur élevé} : erreurs de calcul des taux et montants,
  doublons d'enregistrement, feuilles de maladie incomplètes.
  \item \textbf{Risque de non-conformité} : absence de journalisation et de
  protection structurée des données médicales, données à caractère personnel
  sensibles au sens de la réglementation.
  \item \textbf{Vulnérabilité liée à l'obsolescence} : dépendance à des supports
  non sécurisés et non sauvegardés, risque de perte définitive d'information.
\end{itemize}

\subsubsection{Conséquences de l'inaction}
Le maintien de l'existant exposerait l'organisme à une dégradation continue de la
qualité de service (allongement des délais de remboursement), à un risque
juridique croissant sur la protection des données médicales, et à une incapacité
à passer à l'échelle face à l'augmentation du nombre d'assurés.

\subsection{Objectifs Métiers et Critères de Succès}

Les objectifs ci-dessous sont \textbf{stratégiques, mesurables et datés}. Ils ne
constituent pas des fonctionnalités mais des cibles de performance (KPI) servant à
évaluer la réussite du projet.

\begin{center}
\begin{longtable}{|L{3.0cm}|L{6.0cm}|L{4.6cm}|}
\hline
\rowcolor{anbleu!12}\textbf{Type d'objectif} & \textbf{Exigence structurelle} & \textbf{KPI / Cible datée}\\\hline
\endhead
Objectifs Financiers &
Maîtrise du budget de mise en place (CAPEX) et des coûts opérationnels (OPEX) ;
modèle d'hébergement à coût marginal nul en phase pilote. &
CAPEX projet $\leq$ budget pédagogique alloué ; OPEX d'hébergement $\approx$ 0~FCFA
(offres \emph{free tier} Vercel + base mutualisée) la première année ; ROI mesuré
par la réduction du temps de traitement (cf. ci-dessous) sous 12 mois.\\\hline
Objectifs Opérationnels &
Réduction du temps de traitement, du taux d'erreur et des sollicitations support. &
Réduction $\geq$ 70~\% du délai de constitution d'un dossier de remboursement ;
taux d'erreur de calcul des montants ramené à 0~\% (calcul automatisé) ;
réduction $\geq$ 50~\% des appels support à 6 mois après mise en production.\\\hline
Objectifs Qualitatifs &
Augmentation de la satisfaction des utilisateurs et conformité réglementaire. &
Score de satisfaction (UAT) $\geq$ 4/5 ; conformité aux exigences de protection
des données à caractère personnel applicables au Cameroun dès la mise en
production ; couverture de tests du service métier $\geq$ 70~\% d'instructions.\\\hline
\end{longtable}
\end{center}

\subsection{Glossaire et Ontologie du Projet}

Le glossaire ci-dessous unifie la nomenclature métier et technique afin de
prévenir toute barrière sémantique entre experts métiers et ingénieurs, et de
garantir que chaque entité logicielle corresponde rigoureusement à sa réalité
fonctionnelle.

\begin{center}
\begin{longtable}{|L{3.4cm}|L{11.0cm}|}
\hline
\rowcolor{anbleu!12}\textbf{Terme / Acronyme} & \textbf{Définition}\\\hline
\endhead
Assuré (Patient) & Personne affiliée à l'organisme, identifiée par un numéro de
sécurité sociale (NSS) unique à 15 chiffres. Rôle applicatif : \texttt{PATIENT}.\\\hline
Médecin & Professionnel de santé, généraliste ou spécialiste, identifié par un
numéro RPPS unique à 11 chiffres. Rôle : \texttt{MEDECIN}.\\\hline
Assureur & Agent de sécurité sociale, gestionnaire administratif et financier de
l'application. Rôle : \texttt{ASSUREUR}.\\\hline
Administrateur & Gestionnaire technique du système, accès complet. Rôle : \texttt{ADMIN}.\\\hline
Médecin traitant & Médecin généraliste référent rattaché à un assuré.\\\hline
Consultation & Acte médical enregistré (généraliste ou spécialiste), point de
départ du parcours de soins.\\\hline
Prescription & Ordonnance issue d'une consultation : médicament (\emph{MEDICAMENT})
ou orientation vers un spécialiste (\emph{CONSULTATION\_SPECIALISTE}).\\\hline
Feuille de maladie (Medical Record) & Document médical généré automatiquement à la
création d'une consultation et servant de base au remboursement.\\\hline
Remboursement (Reimbursement) & Demande de prise en charge financière, calculée à
100~\% (généraliste) ou 80~\% (spécialiste), suivant un cycle de vie d'états.\\\hline
Justificatif & Document PDF attestant d'un remboursement.\\\hline
NSS & Numéro de Sécurité Sociale (15 chiffres).\\\hline
RPPS & Répertoire Partagé des Professionnels de Santé (numéro à 11 chiffres).\\\hline
RBAC & \emph{Role-Based Access Control} — contrôle d'accès basé sur les rôles.\\\hline
JWT & \emph{JSON Web Token} — jeton d'authentification signé.\\\hline
MFA & \emph{Multi-Factor Authentication} — authentification multifacteur.\\\hline
SLA & \emph{Service Level Agreement} — accord de niveau de service.\\\hline
PRA / PCA & Plan de Reprise / Continuité d'Activité.\\\hline
RTO / RPO & \emph{Recovery Time / Point Objective} — temps et perte de données maximaux tolérés.\\\hline
CRUD & \emph{Create, Read, Update, Delete} — opérations sur les données.\\\hline
API REST & Interface de programmation respectant le style architectural REST.\\\hline
FCFA & Franc CFA, devise des montants.\\\hline
UAT & \emph{User Acceptance Test} — test / recette d'acceptation utilisateur.\\\hline
RTM & \emph{Requirements Traceability Matrix} — matrice de traçabilité des exigences.\\\hline
\end{longtable}
\end{center}

%==============================================================================
%  SECTION 2
%==============================================================================
\section{Périmètre et Frontières du Système}

Afin de prévenir la \textbf{dérive du périmètre} (\emph{scope creep}), cette section
établit des frontières étanches au projet.

\subsection{Inclusions Explicites (In Scope)}
Le système prend en charge :
\begin{itemize}[leftmargin=1.4em]
  \item la gestion des comptes et l'authentification des quatre profils
  (Administrateur, Médecin, Assureur, Assuré) ;
  \item l'inscription et la gestion des assurés (NSS, affiliation, emploi,
  médecin traitant) et des médecins (RPPS, spécialité) ;
  \item l'enregistrement des consultations et la création automatique de la
  feuille de maladie associée ;
  \item la prescription de médicaments et l'orientation vers un spécialiste ;
  \item la constitution, le calcul automatique, le suivi du cycle de vie et la
  génération du justificatif PDF des remboursements ;
  \item la consultation des feuilles de maladie et l'historique du parcours de soins ;
  \item la journalisation d'audit des opérations sensibles ;
  \item les notifications par courrier électronique aux assurés ;
  \item un tableau de bord statistique pour l'assureur et l'administrateur.
\end{itemize}

\subsection{Exclusions Explicites (Out of Scope)}
Sont \textbf{formellement écartées} de la présente version du projet, et n'engagent
pas la maîtrise d'œuvre lors de la recette :
\begin{itemize}[leftmargin=1.4em]
  \item la prise de rendez-vous en ligne et la téléconsultation vidéo ;
  \item la facturation comptable et l'intégration à un logiciel de comptabilité tiers ;
  \item le paiement électronique en ligne (Mobile Money, carte bancaire) — seuls
  les modes \emph{Virement} et \emph{Espèces} sont enregistrés à titre déclaratif ;
  \item l'interfaçage avec des systèmes nationaux d'assurance maladie externes ;
  \item l'application mobile native (iOS/Android) — l'accès se fait via navigateur
  web responsive ;
  \item l'authentification multifacteur (MFA) et le verrouillage de compte, prévus
  pour une évolution ultérieure (cf. §8) ;
  \item la gestion multilingue : l'interface est livrée en \textbf{français} uniquement ;
  \item toute intégration avec des pharmacies, laboratoires ou objets connectés.
\end{itemize}

\subsection{Hypothèses, Dépendances et Contraintes}

\textbf{Hypothèses de travail :}
\begin{itemize}[leftmargin=1.4em]
  \item disponibilité d'une connexion Internet sur les postes des utilisateurs
  (l'application est exclusivement en ligne) ;
  \item les utilisateurs disposent d'un navigateur récent (Chrome, Edge, Firefox) ;
  \item les données de référence (assurés, médecins) sont saisies par l'assureur ou
  l'administrateur.
\end{itemize}

\textbf{Dépendances critiques :}
\begin{itemize}[leftmargin=1.4em]
  \item plateforme d'hébergement frontend \textbf{Vercel} ;
  \item serveur d'application \textbf{Tomcat 10.1} exécutant l'archive Spring Boot ;
  \item serveur de base de données \textbf{PostgreSQL} ;
  \item service de messagerie \textbf{SMTP Google (Gmail)} pour les notifications.
\end{itemize}

\textbf{Contraintes imposées :}
\begin{itemize}[leftmargin=1.4em]
  \item \emph{Technologiques} : socle Java~17 / Spring~Boot~3.2.x côté serveur ;
  Next.js~14 / React~18 / TypeScript côté client.
  \item \emph{Budgétaires} : recours prioritaire aux offres gratuites (contexte académique).
  \item \emph{Temporelles} : réalisation complète en \textbf{3 mois} (cf. §10).
  \item \emph{Réglementaires} : conformité aux lois camerounaises sur la
  cybersécurité et la protection des données à caractère personnel (cf. §8 et §11).
\end{itemize}

%==============================================================================
%  SECTION 3
%==============================================================================
\section{Description Générale et Modèle Opérationnel}

\subsection{L'Écosystème et la Perspective du Produit}

\textsc{Assurance Nation} est une \textbf{application web full-stack à architecture
client-serveur découplée}, et non un progiciel autonome. Elle se compose :
\begin{itemize}[leftmargin=1.4em]
  \item d'un \textbf{client web} (Single Page Application Next.js/React) hébergé sur
  Vercel, assurant la présentation et l'export PDF côté navigateur ;
  \item d'une \textbf{API REST stateless} (Spring Boot, exposée en JSON) portant
  l'ensemble de la logique métier et de la sécurité ;
  \item d'une \textbf{base de données relationnelle} PostgreSQL persistant les données ;
  \item d'un \textbf{service de notification} s'appuyant sur le SMTP Gmail.
\end{itemize}

Le diagramme de contexte de haut niveau ci-dessous présente le système comme une
boîte noire centrale et ses flux d'interaction avec les acteurs humains et machines.

\begin{figure}[H]
\centering
\begin{tikzpicture}[
  node distance=1.2cm and 2.4cm,
  acteur/.style={draw,anbleu,thick,rounded corners,align=center,fill=anbleu!6,minimum height=1.0cm,minimum width=2.6cm},
  systeme/.style={draw,anbleu,very thick,fill=anjaune!25,align=center,minimum height=2.0cm,minimum width=3.6cm},
  ext/.style={draw,angris,thick,dashed,align=center,minimum height=1.0cm,minimum width=2.8cm},
  flux/.style={-Stealth,thick,anbleu}]

  \node[systeme] (sys) {\textbf{ASSURANCE}\\\textbf{NATION}\\\small(Système)};
  \node[acteur,left=of sys] (med) {Médecin};
  \node[acteur,above=of sys] (ass) {Assureur};
  \node[acteur,right=of sys] (pat) {Assuré\\(Patient)};
  \node[acteur,below=of sys] (adm) {Administrateur};
  \node[ext,right=3.0cm of pat] (smtp) {Serveur SMTP\\Gmail};

  \draw[flux] (med) -- node[above,sloped,font=\scriptsize]{consultations / prescriptions} (sys);
  \draw[flux] (ass) -- node[right,font=\scriptsize]{remboursements / inscriptions} (sys);
  \draw[flux] (pat) -- node[above,sloped,font=\scriptsize]{consultation dossiers} (sys);
  \draw[flux] (adm) -- node[right,font=\scriptsize]{administration / audit} (sys);
  \draw[flux] (sys.east) .. controls +(2,0) and +(-1,-1) .. node[below,font=\scriptsize]{notifications} (smtp.south west);
\end{tikzpicture}
\caption{Diagramme de contexte de haut niveau du système \textsc{Assurance Nation}.}
\end{figure}

\subsection{Typologie des Utilisateurs et Caractéristiques}

\begin{center}
\begin{longtable}{|L{2.6cm}|L{3.0cm}|L{3.0cm}|L{4.4cm}|}
\hline
\rowcolor{anbleu!12}\textbf{Profil} & \textbf{Expertise technique} & \textbf{Fréquence d'usage} & \textbf{Environnement d'opération}\\\hline
\endhead
Médecin & Moyenne & Quotidienne (en consultation) & Poste de bureau ou ordinateur portable au cabinet ; sessions courtes et répétées.\\\hline
Assureur & Moyenne à élevée & Intensive (toute la journée) & Poste sédentaire, souvent double écran ; traitement de volumes de dossiers.\\\hline
Assuré (Patient) & Faible & Occasionnelle & Domicile, navigateur grand public ; doit bénéficier d'une interface simple et guidée.\\\hline
Administrateur & Élevée & Ponctuelle & Poste technique ; supervision, gestion des comptes et audit.\\\hline
\end{longtable}
\end{center}

Cette classification fonde les exigences d'ergonomie et d'accessibilité (cf. §5) :
priorité à la simplicité pour l'assuré, à la densité d'information pour l'assureur.

%==============================================================================
%  SECTION 4
%==============================================================================
\section{Exigences Fonctionnelles (Spécifications Métier)}

Les exigences fonctionnelles définissent le \textbf{« Quoi »}. Conformément à la
norme \textbf{ISO/IEC/IEEE~29148}, chaque exigence est atomique, vérifiable et
documentée selon un schéma d'attributs strict (identifiant unique, intitulé,
description impérative, priorité MoSCoW, complexité, critères d'acceptation, méthode
de vérification).

\subsection{Anatomie d'une Exigence Normée — Domaine Authentification \& Sécurité}

\exig{EF-AUTH-01}{Inscription d'un utilisateur}
{Le système doit permettre l'enregistrement d'un utilisateur (Médecin, Assuré ou
Assureur) avec validation des champs obligatoires et unicité de l'adresse
électronique.}
{Must (Indispensable)}{Moyenne}
{Un compte est créé ; l'e-mail en doublon est rejeté avec un message explicite ;
le mot de passe respecte la politique (8+ caractères, 1 majuscule, 1 chiffre, 1
caractère spécial) ; un couple jeton d'accès + jeton de rafraîchissement est retourné.}
{Test automatisé + démonstration}

\exig{EF-AUTH-02}{Authentification (connexion)}
{Le système doit authentifier un utilisateur par e-mail et mot de passe et délivrer
un jeton JWT signé.}
{Must}{Faible}
{Des identifiants valides ouvrent une session et renvoient un jeton d'accès
(durée 15~min) et un jeton de rafraîchissement (durée 7~jours) ; des identifiants
invalides retournent une erreur 401 sans divulgation.}
{Test automatisé}

\exig{EF-AUTH-03}{Rafraîchissement et déconnexion}
{Le système doit renouveler un jeton d'accès expiré via le jeton de rafraîchissement
et révoquer ce dernier à la déconnexion.}
{Must}{Faible}
{Un jeton de rafraîchissement valide produit un nouveau jeton d'accès ; après
déconnexion, le jeton de rafraîchissement est révoqué et inutilisable.}
{Test automatisé}

\subsection{Domaine Gestion des Utilisateurs}

\exig{EF-USR-01}{Gestion des assurés}
{Le système doit permettre de créer, rechercher (par nom ou NSS), consulter, mettre
à jour et supprimer logiquement un assuré.}
{Must}{Moyenne}
{Le NSS (15 chiffres) est unique et validé ; la recherche par NSS retourne l'assuré ;
la suppression est logique (\emph{soft delete}) et exclut l'assuré des listes.}
{Test automatisé + inspection}

\exig{EF-USR-02}{Gestion des médecins}
{Le système doit permettre de créer, rechercher (par nom, RPPS ou spécialité),
consulter, mettre à jour et supprimer logiquement un médecin.}
{Must}{Moyenne}
{Le RPPS (11 chiffres) est unique ; un médecin spécialiste doit posséder un libellé
de spécialité ; la suppression est logique.}
{Test automatisé}

\exig{EF-USR-03}{Affectation du médecin traitant}
{Le système doit permettre à l'assureur ou l'administrateur d'affecter ou de modifier
le médecin traitant d'un assuré.}
{Should (Important)}{Faible}
{Le médecin traitant est associé à l'assuré et restitué dans sa fiche.}
{Démonstration}

\exig{EF-USR-04}{Gestion du profil et du mot de passe}
{Le système doit permettre à un utilisateur de consulter et modifier son profil et
de changer son mot de passe après vérification du mot de passe courant.}
{Must}{Faible}
{La modification est restreinte au propriétaire ou à l'administrateur ; un mauvais
mot de passe courant bloque le changement.}
{Test automatisé}

\subsection{Domaine Consultations et Prescriptions}

\exig{EF-CONS-01}{Création d'une consultation}
{Le système doit permettre à un médecin de créer une consultation (généraliste ou
spécialiste) pour un assuré et de générer automatiquement la feuille de maladie associée.}
{Must}{Élevée}
{Une consultation est enregistrée avec son type ; une feuille de maladie liée est
créée avec \texttt{estRemboursee=false} ; une notification est envoyée à l'assuré.}
{Test automatisé + démonstration}

\exig{EF-CONS-02}{Consultation et annulation}
{Le système doit permettre de lister/consulter les consultations selon le rôle, et au
médecin d'annuler (suppression logique) une consultation.}
{Must}{Moyenne}
{Le médecin ne voit que ses consultations, le patient les siennes, l'assureur/admin
toutes ; l'annulation est logique et notifiée.}
{Test automatisé}

\exig{EF-PRESC-01}{Prescription de médicament}
{Le système doit permettre au médecin d'ajouter à une consultation une prescription
de médicament (nom, posologie, durée 1--365 jours, notes).}
{Must}{Moyenne}
{La prescription est rattachée à la consultation ; le médicament et la durée sont
obligatoires ; l'assuré est notifié.}
{Test automatisé}

\exig{EF-PRESC-02}{Orientation vers un spécialiste}
{Le système doit permettre au médecin de prescrire une consultation chez un
spécialiste (motif, priorité, code de référence, spécialiste optionnel).}
{Must}{Moyenne}
{La prescription de type \emph{CONSULTATION\_SPECIALISTE} est créée avec un code de
référence ; le motif est obligatoire.}
{Test automatisé}

\subsection{Domaine Feuilles de Maladie et Remboursements}

\exig{EF-FM-01}{Consultation des feuilles de maladie}
{Le système doit permettre de lister et consulter les feuilles de maladie d'un assuré
avec filtres (période, statut de remboursement).}
{Must}{Faible}
{Les filtres par date et statut s'appliquent ; une feuille déjà remboursée n'est plus
modifiable.}
{Test automatisé}

\exig{EF-RB-01}{Création d'un remboursement}
{Le système doit permettre à l'assureur de créer un remboursement à partir d'une
feuille de maladie, avec calcul automatique du montant (100~\% généraliste, 80~\%
spécialiste) et génération du justificatif PDF.}
{Must}{Élevée}
{Le remboursement est créé au statut \emph{PENDING} avec un numéro unique ; le montant
remboursé est correctement calculé ; un PDF justificatif est disponible ; une feuille
déjà remboursée est refusée.}
{Test automatisé + démonstration}

\exig{EF-RB-02}{Cycle de vie du remboursement (workflow)}
{Le système doit gérer les transitions d'état du remboursement et notifier l'assuré à
chaque étape.}
{Must}{Élevée}
{Transitions autorisées : PENDING$\to$APPROVED, APPROVED$\to$PAID, PENDING$\to$REJECTED
(avec motif) ; toute transition illégale est rejetée ; un courriel est envoyé à l'assuré.}
{Test automatisé}

\exig{EF-RB-03}{Téléchargement du justificatif}
{Le système doit permettre le téléchargement du justificatif PDF d'un remboursement.}
{Should}{Faible}
{Le document PDF est restitué avec le type MIME \texttt{application/pdf}.}
{Démonstration}

\exig{EF-RB-04}{Tableau de bord des remboursements}
{Le système doit fournir un tableau de bord statistique (compteurs par statut, montants,
tendance mensuelle, répartition par spécialité).}
{Should}{Moyenne}
{Les compteurs et montants correspondent aux données ; les graphiques s'affichent.}
{Démonstration + analyse}

\subsection{Domaine Audit et Notifications}

\exig{EF-AUD-01}{Journal d'audit inaltérable}
{Le système doit historiser toute opération sensible (création, modification,
suppression) sur les entités Utilisateur, Consultation, Prescription et Remboursement.}
{Must}{Moyenne}
{Chaque opération produit une entrée d'audit (type d'entité, identifiant, action, ancien
et nouveau état, horodatage, adresse IP, acteur) ; le journal est consultable par
l'administrateur et non modifiable.}
{Inspection + test automatisé}

\exig{EF-NOT-01}{Notifications par courrier électronique}
{Le système doit notifier l'assuré lors des événements clés (consultation, prescription,
étapes de remboursement) via le SMTP Gmail, de façon asynchrone.}
{Should}{Moyenne}
{Un courriel est émis pour chaque événement ; un échec d'envoi ne bloque pas l'opération
métier.}
{Démonstration}

\subsection{Organisation par Cas d'Utilisation et Workflows}

\subsubsection{Cas d'utilisation — Effectuer un remboursement (UC-RB)}
\begin{description}[leftmargin=2.2cm,style=nextline]
  \item[Acteur principal] Assureur (ou Administrateur).
  \item[Préconditions] L'assureur est authentifié ; une feuille de maladie non encore
  remboursée existe pour l'assuré.
  \item[Scénario nominal]
  \begin{enumerate}[leftmargin=1.6em]
    \item L'assureur recherche l'assuré par son NSS.
    \item Il sélectionne la feuille de maladie concernée.
    \item Il saisit le montant total et le mode de paiement (Virement/Espèces).
    \item Le système calcule le montant remboursé selon le type de consultation.
    \item Le système crée le remboursement au statut \emph{PENDING}, génère le PDF et
    notifie l'assuré.
    \item L'assureur approuve, puis marque comme payé.
  \end{enumerate}
  \item[Scénarios alternatifs]
  \begin{itemize}[leftmargin=1.6em]
    \item A1 : la feuille est déjà remboursée $\Rightarrow$ le système refuse la création.
    \item A2 : montant invalide ($\leq 0$) $\Rightarrow$ message de validation.
    \item A3 : rejet de la demande $\Rightarrow$ statut \emph{REJECTED} avec motif obligatoire.
  \end{itemize}
  \item[Post-conditions] Le remboursement est persisté dans son état final ; la feuille
  de maladie est marquée remboursée ; les événements sont audités et notifiés.
\end{description}

\subsubsection{Modélisation du workflow de remboursement (BPM)}
\begin{figure}[H]
\centering
\begin{tikzpicture}[node distance=2.2cm,
  etat/.style={draw,anbleu,thick,rounded corners,fill=anbleu!8,minimum width=2.2cm,minimum height=0.9cm,align=center},
  fin/.style={draw,thick,rounded corners,minimum width=2.0cm,minimum height=0.9cm,align=center},
  fl/.style={-Stealth,thick}]
  \node[etat] (p) {PENDING};
  \node[etat,right=of p] (a) {APPROVED};
  \node[fin,fill=green!12,right=of a] (pa) {PAID};
  \node[fin,fill=red!12,below=1.3cm of p] (r) {REJECTED};
  \draw[fl] (p) -- node[above,font=\scriptsize]{approuver} (a);
  \draw[fl] (a) -- node[above,font=\scriptsize]{marquer payé} (pa);
  \draw[fl] (p) -- node[left,font=\scriptsize]{rejeter (motif)} (r);
\end{tikzpicture}
\caption{Workflow d'approbation du remboursement : événements déclencheurs et
notifications associées à chaque transition.}
\end{figure}

\subsection{Exigences sur le Modèle et le Cycle de Vie des Données}

Les traitements reposent sur un modèle de données relationnel. Les entités
principales, leurs attributs clés et leurs contraintes d'intégrité sont formalisés
ci-dessous.

\begin{center}
\begin{longtable}{|L{3.0cm}|L{7.0cm}|L{4.0cm}|}
\hline
\rowcolor{anbleu!12}\textbf{Entité} & \textbf{Attributs principaux} & \textbf{Contraintes \& relations}\\\hline
\endhead
\texttt{User} (abstrait) & id (UUID), email, passwordHash, nom, prénom, dateNaissance,
adresse, téléphone, sexe, photoUrl, horodatages d'audit & email \textbf{unique, non nul} ;
N--N avec \texttt{Role} ; sous-types Médecin / Assuré / Assureur.\\\hline
\texttt{Medecin} & numéroRPPS, spécialité (GENERALISTE/SPECIALISTE), libellé, estAssure &
RPPS \textbf{unique} ; 1--N \texttt{Consultation} ; 1--N assurés suivis.\\\hline
\texttt{Assure} & numSécuritéSociale, dateAffiliation, emploi, estActif &
NSS \textbf{unique} ; N--1 médecin traitant ; 1--N feuilles de maladie.\\\hline
\texttt{Consultation} & dateConsultation, type, diagnostic, motif, notes &
N--1 Assuré, N--1 Médecin ; 1--N Prescriptions ; 1--1 Feuille de maladie.\\\hline
\texttt{Prescription} (abstrait) & type & N--1 Consultation ; sous-types
\texttt{PrescriptionMedicament} et \texttt{PrescriptionConsultation}.\\\hline
\texttt{MedicalRecord} & date, nomMaladie, estRemboursee, montantRemboursé &
1--1 Consultation \textbf{unique} ; N--1 Assuré et Médecin.\\\hline
\texttt{Reimbursement} & numRemboursement, montantTotal, taux, montantRemboursé,
modePaiement, status & numéro \textbf{unique} ; 1--1 Feuille de maladie ; statut $\in$
\{PENDING, APPROVED, PAID, REJECTED\}.\\\hline
\texttt{Role / Permission} & nom, description & N--N ; jeu de rôles ADMIN, MEDECIN,
ASSUREUR, PATIENT.\\\hline
\texttt{RefreshToken} & token, expiryDate, revoked & N--1 User.\\\hline
\texttt{AuditLog} & entityType, entityId, action, oldValue, newValue, timestamp, ip &
N--1 User (acteur) ; \textbf{inaltérable}.\\\hline
\end{longtable}
\end{center}

\textbf{Règles de cycle de vie (CRUD) :} la suppression des entités sensibles est
\textbf{logique} (\emph{soft delete} via \texttt{deleted\_at}) afin de préserver la
piste d'audit ; aucune donnée n'est physiquement détruite. Les feuilles de maladie
remboursées sont en lecture seule. L'intégrité référentielle est garantie par clés
étrangères et contraintes d'unicité (email, NSS, RPPS, numéro de remboursement).

%==============================================================================
%  SECTION 5
%==============================================================================
\section{Exigences Non-Fonctionnelles (Attributs de Qualité)}

Les exigences non-fonctionnelles dictent le \textbf{« Comment »} et garantissent que le
logiciel survivra à son environnement d'exploitation.

\subsection{Performance, Tolérance et Scalabilité}
\begin{itemize}[leftmargin=1.4em]
  \item \textbf{Temps de réponse} : opérations courantes (lecture, création unitaire)
  $\leq$ 1,5~s au 95\textsuperscript{e} percentile ; requêtes complexes (tableau de
  bord, recherches filtrées) $\leq$ 3~s.
  \item \textbf{Pagination} : toutes les listes sont paginées (taille par défaut 20,
  maximum 100) pour borner la charge.
  \item \textbf{Concurrence nominale} : le système doit supporter au moins
  \textbf{100 utilisateurs concurrents} sans dégradation perceptible.
  \item \textbf{Volumétrie projetée} : capacité d'absorber $\geq$ 100\,000 assurés,
  500\,000 consultations et 500\,000 remboursements sur une projection de 3 ans.
  \item \textbf{Statelessness} : l'API étant sans état (JWT), elle peut être mise à
  l'échelle horizontalement.
\end{itemize}

\subsection{Fiabilité, Disponibilité et Maintien en Condition Opérationnelle}
\begin{itemize}[leftmargin=1.4em]
  \item \textbf{Disponibilité (SLA)} : taux de disponibilité cible de \textbf{99,5~\%}
  mensuel en phase pilote (objectif 99,9~\% en production cible).
  \item \textbf{Fenêtre de maintenance} : interruptions planifiées $\leq$ 4~h/mois,
  notifiées 48~h à l'avance, hors heures ouvrées.
  \item \textbf{Mode dégradé} : un échec du service de notification (SMTP) ne doit pas
  interrompre les opérations métier (envoi asynchrone, erreur journalisée).
  \item \textbf{Tolérance aux pannes} : la couche d'accès aux données et les transactions
  garantissent la cohérence ; un point de santé (\texttt{/health}) expose l'état de l'API,
  de la base et de l'espace disque.
\end{itemize}

\subsection{Ergonomie, Expérience Utilisateur et Accessibilité}
\begin{itemize}[leftmargin=1.4em]
  \item \textbf{Charte graphique} : respect des couleurs institutionnelles
  (bleu primaire \texttt{\#0066CC}, bleu marine \texttt{\#003366}, accent vert
  \texttt{\#00AA44}) et de la typographie de marque.
  \item \textbf{Conception centrée utilisateur} : parcours simplifiés pour l'assuré,
  vues denses pour l'assureur (cf. §3.2).
  \item \textbf{Responsive design} : adaptation mobile (320~px), tablette et bureau
  (1024~px+).
  \item \textbf{Thème} : modes clair et sombre, préférence persistée.
  \item \textbf{Accessibilité} : conformité visée au niveau \textbf{WCAG~2.1~AA}
  (contrastes, navigation clavier, libellés de formulaires, attributs ARIA).
\end{itemize}

\subsection{Maintenabilité, Portabilité et Compatibilité}
\begin{itemize}[leftmargin=1.4em]
  \item \textbf{Navigateurs supportés} : Chrome, Edge, Firefox (deux dernières versions
  majeures).
  \item \textbf{Frameworks pérennes} : Spring Boot 3.2.x (LTS Java 17), Next.js 14
  (React 18), TypeScript, choisis pour leur support à long terme.
  \item \textbf{Standards de codage} : conventions Java et ESLint/TypeScript ; revue de
  code systématique.
  \item \textbf{Documentation technique} : API documentée via OpenAPI/Swagger ;
  documentation du modèle de données et des diagrammes UML livrée.
  \item \textbf{Portabilité} : application conteneurisable (Docker) ; persistance via JPA
  abstrayant le SGBD relationnel.
  \item \textbf{Testabilité} : couverture de tests du service métier $\geq$ 70~\%
  d'instructions.
\end{itemize}

%==============================================================================
%  SECTION 7  (la 6 est volontairement omise)
%==============================================================================
\setcounter{section}{6}
\section{Exigences Relatives à l'Hébergement et aux Architectures Cloud / SaaS}

\subsection{Modèle de Déploiement et Ségrégation des Données}

Le système adopte un \textbf{modèle de déploiement en cloud public}, avec séparation
nette des responsabilités :
\begin{itemize}[leftmargin=1.4em]
  \item \textbf{Frontend} hébergé sur \textbf{Vercel} (rendu Next.js, prérendu de
  contenu statique, CDN) ;
  \item \textbf{Backend} packagé en archive exécutable Spring Boot et servi par
  \textbf{Tomcat 10.1} sur la JVM ;
  \item \textbf{Base de données} \textbf{PostgreSQL}, instance \texttt{assurance\_nation} ;
  \item \textbf{Notifications} via le serveur \textbf{SMTP Google (Gmail)}.
\end{itemize}

Les flux et les ports normalisés de l'architecture de déploiement sont :
\begin{enumerate}[leftmargin=1.6em]
  \item Navigateur $\to$ Frontend Vercel : \textbf{HTTPS / port 443} ;
  \item Navigateur $\to$ API Backend : \textbf{REST / HTTPS / JSON / port 8083} ;
  \item Backend $\to$ Base de données : \textbf{JPA / JDBC / TCP / port 5432} ;
  \item Backend $\to$ SMTP Gmail : \textbf{TLS / STARTTLS / port 587}.
\end{enumerate}

\textbf{Ségrégation et étanchéité des données :} l'instance de base de données est
\textbf{dédiée} à l'application (\emph{single-tenant} logique). Dans l'hypothèse d'un
hébergement mutualisé, l'isolation est garantie par un schéma propre, des comptes de
service à privilèges minimaux et le chiffrement des connexions (JDBC sur TLS), afin
d'éviter toute porosité d'informations entre clients.

\subsection{Schéma d'Architecture de Déploiement}

Le schéma ci-dessous formalise l'architecture cible (postes clients, hébergement
frontend, serveur d'application, base de données et service de messagerie) ainsi que
les protocoles et ports d'interconnexion.

\begin{figure}[H]
\centering
\resizebox{\linewidth}{!}{%
\begin{tikzpicture}[font=\scriptsize,
  comp/.style={draw,align=center,minimum height=0.9cm,minimum width=2.8cm,fill=white,inner sep=3pt},
  dbc/.style={draw,cylinder,shape border rotate=90,aspect=0.28,align=center,minimum width=2.8cm,minimum height=1.3cm,fill=white},
  cont/.style={draw,thick,rounded corners=2pt,fill=anbleu!4},
  lbl/.style={font=\scriptsize\bfseries\color{anbleu}},
  fl/.style={-Stealth,thick,anbleu},
  dfl/.style={-Stealth,thick,dotted},
  flab/.style={font=\scriptsize,align=center,fill=white,inner sep=1.5pt}]

  % --- Poste client ---
  \node[comp,minimum width=6.8cm] (nav) at (6,9.4) {Navigateur Web : Chrome, Edge, Firefox};
  \node[comp] (react) at (4.2,8.0) {React State /\\LocalStorage};
  \node[comp] (pdf)   at (7.8,8.0) {jsPDF / html2canvas\\(Export PDF)};
  \begin{scope}[on background layer]
    \node[cont,fit=(nav)(react)(pdf),inner sep=9pt] (client) {};
  \end{scope}
  \node[lbl,anchor=south west,yshift=1pt] at (client.north west) {Poste Client};

  % --- Cloud Vercel ---
  \node[comp] (next)   at (1.8,4.6) {Next.js App Server /\\Node.js Engine};
  \node[comp] (assets) at (6.5,4.6) {Static HTML/CSS/JS\\Assets};
  \draw[dfl] (next) -- node[above,font=\scriptsize,fill=white,inner sep=1pt]{prérendu} (assets);
  \begin{scope}[on background layer]
    \node[cont,fit=(next)(assets),inner sep=9pt] (vercel) {};
  \end{scope}
  \node[lbl,anchor=south west,yshift=1pt] at (vercel.north west) {Cloud Vercel (Hébergement Frontend)};

  % --- Serveur d'application ---
  \node[comp,minimum width=3.8cm] (spring) at (11.9,5.5) {SpringBootApplication.jar\\Tomcat 10.1};
  \node[comp,minimum width=3.8cm] (jvm)    at (11.9,4.1) {Java Virtual Machine (JVM)};
  \draw[fl] (spring) -- (jvm);
  \begin{scope}[on background layer]
    \node[cont,fit=(spring)(jvm),inner sep=9pt] (appsrv) {};
  \end{scope}
  \node[lbl,anchor=south,yshift=1pt] at (appsrv.north) {Serveur d'application};

  % --- Serveur de base de données ---
  \node[comp,minimum width=3.2cm] (pg)  at (5.3,1.4) {PostgreSQL\\Database Server};
  \node[dbc] (dbi) at (5.3,-0.9) {Database Instance:\\\texttt{assurance\_nation}};
  \draw[fl] (pg) -- node[right]{DBMS Node} (dbi);
  \begin{scope}[on background layer]
    \node[cont,fit=(pg)(dbi),inner sep=9pt] (dbsrv) {};
  \end{scope}
  \node[lbl,anchor=south west,yshift=1pt] at (dbsrv.north west) {Serveur de base de données};

  % --- Serveur SMTP ---
  \node[comp,minimum width=3.2cm] (smtp) at (12.5,-0.3) {smtp.gmail.com\\Port 587};
  \begin{scope}[on background layer]
    \node[cont,fit=(smtp),inner sep=9pt] (smtpsrv) {};
  \end{scope}
  \node[lbl,anchor=south west,yshift=1pt] at (smtpsrv.north west) {Serveur SMTP Google (Gmail)};

  % --- Flux ---
  \draw[fl] (3.4,0|-client.south) -- node[flab,pos=0.5]{1. HTTPS\\Port 443}
            (3.4,0|-vercel.north);
  \draw[fl] (8.6,0|-client.south) -- node[flab,pos=0.45]{2. REST / HTTPS\\JSON / Port 8083}
            (appsrv.north west);
  \draw[fl] (appsrv.south west) -- node[flab,pos=0.5]{3. JPA / JDBC\\TCP / Port 5432}
            (dbsrv.east);
  \draw[fl] ([xshift=0.6cm]appsrv.south) -- node[flab,pos=0.55]{4. TLS / STARTTLS\\Port 587}
            (smtpsrv.north);
\end{tikzpicture}}
\caption{Diagramme de déploiement de la plateforme \textsc{Assurance Nation} :
navigateur web, hébergement frontend Vercel, serveur d'application Spring~Boot /
Tomcat~10.1, base de données PostgreSQL (\texttt{assurance\_nation}) et serveur SMTP
Gmail, avec les protocoles et ports normalisés.}
\end{figure}

\subsection{Plans de Continuité, Sauvegardes et Réversibilité}

\textbf{PRA / PCA :}
\begin{itemize}[leftmargin=1.4em]
  \item \textbf{RTO} (temps de récupération maximal) : $\leq$ 4~heures pour le service
  applicatif ;
  \item \textbf{RPO} (perte de données maximale tolérée) : $\leq$ 24~heures (aligné sur
  la fréquence de sauvegarde) ;
  \item procédure de redéploiement documentée (image conteneur + restauration de base)
  permettant une reprise sur incident majeur.
\end{itemize}

\textbf{Politique de sauvegarde :}
\begin{itemize}[leftmargin=1.4em]
  \item sauvegarde complète \textbf{quotidienne} de la base \texttt{assurance\_nation} ;
  \item \textbf{externalisation sécurisée} des supports (stockage chiffré distinct du
  serveur de production) ;
  \item \textbf{rétention} de 30 jours glissants minimum, avec test de restauration
  périodique.
\end{itemize}

\textbf{Réversibilité :} en fin de contrat, le client peut récupérer
\textbf{l'intégralité de ses données} dans un format interopérable et non chiffré
(export SQL et/ou CSV/JSON), accompagné du schéma de données et de la documentation,
sans aucune rétention abusive de la part de l'éditeur. Les justificatifs PDF sont
exportés tels quels.

%==============================================================================
%  SECTION 8
%==============================================================================
\section{Cybersécurité et Conformité Réglementaire}

\subsection{Ingénierie de Sécurité (Security by Design)}
La sécurité est intégrée dès la conception :
\begin{itemize}[leftmargin=1.4em]
  \item \textbf{Protection contre les vulnérabilités classiques} : prévention des
  injections SQL par l'usage de requêtes paramétrées (JPA/ORM) ; protection contre les
  failles XSS par l'échappement des sorties côté client et la validation des entrées
  (Zod côté frontend, Bean Validation côté backend).
  \item \textbf{Chiffrement en transit} : l'ensemble des communications externes
  s'effectue en HTTPS/TLS (ports 443 et 8083) ; les courriels via STARTTLS (587) ; la
  connexion base via JDBC sur TLS.
  \item \textbf{Chiffrement au repos} : chiffrement des volumes de stockage et des
  sauvegardes ; les données sensibles bénéficient des mécanismes de chiffrement du SGBD.
  \item \textbf{Stockage des mots de passe} : algorithme de hachage robuste et salé
  \textbf{BCrypt} (facteur de coût 12). Aucun mot de passe n'est stocké en clair.
\end{itemize}

\subsection{Gestion des Habilitations, Traçabilité et Audit}
\begin{itemize}[leftmargin=1.4em]
  \item \textbf{Contrôle d'accès basé sur les rôles (RBAC)} : quatre rôles (ADMIN,
  MEDECIN, ASSUREUR, PATIENT) ; application stricte du \textbf{principe du moindre
  privilège} via la sécurité au niveau des méthodes (\texttt{@PreAuthorize}) et le
  contrôle de propriété (un médecin n'accède qu'à ses consultations, un assuré qu'à
  ses propres dossiers).
  \item \textbf{Authentification} : jetons JWT signés (HMAC-SHA256), session sans état,
  jetons d'accès de courte durée (15~min) et jetons de rafraîchissement révocables.
  \item \textbf{Évolutions prévues} : authentification multifacteur (MFA) et
  verrouillage de compte après tentatives infructueuses (cf. §2.2 — hors périmètre de
  la version courante, recommandés pour la mise en production réelle).
  \item \textbf{Journaux d'audit inaltérables} : chaque accès ou modification de donnée
  sensible est historisé (qui, quoi, quand, depuis quelle adresse IP, ancien/nouvel
  état), conformément à l'exigence EF-AUD-01.
\end{itemize}

\subsection{Conformité Réglementaire (Protection des Données)}
Le traitement de données de santé, particulièrement sensibles, impose :
\begin{itemize}[leftmargin=1.4em]
  \item le respect de la \textbf{loi camerounaise n\textsuperscript{o}~2024/017 du
  23 décembre 2024} régissant la protection des données à caractère personnel, et de
  la \textbf{loi n\textsuperscript{o}~2010/012 du 21 décembre 2010} relative à la
  cybersécurité et à la cybercriminalité au Cameroun, sous l'égide de l'\textbf{ANTIC} ;
  \item la prise en compte des principes du \textbf{RGPD} (UE) comme état de l'art :
  minimisation, finalité, consentement, droit d'accès et de rectification, durée de
  conservation limitée ;
  \item la confidentialité du secret médical ; l'accès aux données de santé est
  strictement restreint par le RBAC et tracé par l'audit.
\end{itemize}

%==============================================================================
%  SECTION 10  (la 9 est volontairement omise)
%==============================================================================
\setcounter{section}{9}
\section{Méthodologie, Gestion de Projet et Assurance Qualité}

\subsection{Cadre Méthodologique et Cycle de Développement}
Le projet adopte une \textbf{approche Agile itérative} adaptée au contexte académique
(itérations hebdomadaires, démonstrations régulières, gestion d'un \emph{backlog}
d'exigences), tout en s'appuyant sur des \textbf{jalons documentaires} hérités d'un
cycle en V pour les phases d'analyse et de conception (cahier d'analyse UML approuvé
avant développement).

\textbf{Rôles :}
\begin{itemize}[leftmargin=1.4em]
  \item \textbf{Chef de projet / Product Owner} : TAGNE TEGUEO KUATE FREEDY-JULIO —
  priorisation du backlog, interface avec l'encadrement.
  \item \textbf{Encadrement (maîtrise d'ouvrage)} : Dr. Anne Marie CHANA,
  Dr. Jaures Styve KAMENI — validation des jalons.
  \item \textbf{Équipe de réalisation} : membres du Groupe 9, répartis en sous-équipes
  Analyse/Conception, Backend, Frontend et Tests.
\end{itemize}

\subsection{Stratégie de Test et Plan de Validation}
\begin{itemize}[leftmargin=1.4em]
  \item \textbf{Tests unitaires automatisés} : JUnit~5 (backend), Vitest (frontend) —
  cible $\geq$ 70~\% d'instructions sur la couche service.
  \item \textbf{Tests d'intégration} : Spring Boot Test + Testcontainers ; REST Assured
  pour les contrats d'API.
  \item \textbf{Tests de bout en bout (E2E)} : Cypress (parcours connexion, création de
  consultation, demande de remboursement).
  \item \textbf{Tests de charge} : simulation d'au moins 100 utilisateurs concurrents
  (validation de l'exigence de performance).
  \item \textbf{Recette / UAT} : validation par l'encadrement sur les critères
  d'acceptation des exigences fonctionnelles.
\end{itemize}

\textbf{Classification des anomalies et délais de correction :}
\begin{center}
\begin{longtable}{|L{2.6cm}|L{7.4cm}|L{3.4cm}|}
\hline
\rowcolor{anbleu!12}\textbf{Gravité} & \textbf{Définition} & \textbf{Délai de correction}\\\hline
\endhead
Bloquante & Empêche l'usage d'une fonction majeure ; aucun contournement. & 24~h ; bloque la mise en production.\\\hline
Majeure & Dysfonctionnement important avec contournement possible. & 72~h.\\\hline
Mineure & Défaut cosmétique ou de confort sans impact métier. & Prochaine itération.\\\hline
\end{longtable}
\end{center}

\subsection{Accompagnement, Déploiement et Formation}
Le plan d'accompagnement au changement comprend la production des \textbf{manuels
utilisateurs} (par profil), de la \textbf{documentation technique} (API OpenAPI,
modèle de données, guide de déploiement) et de sessions de prise en main.

\subsection{Planning Prévisionnel — Diagramme de Gantt (3 mois)}
Le diagramme ci-dessous présente le planning réel du projet, de la conception
(diagrammes UML) jusqu'au code et au déploiement, sur un horizon de \textbf{3 mois}
(12 semaines).

\begin{figure}[H]
\centering
\resizebox{\linewidth}{!}{%
\begin{ganttchart}[
  hgrid, vgrid,
  bar/.append style={fill=anbleu!55,draw=anbleu},
  bar height=0.55,
  group/.append style={fill=anjaune!70,draw=anjaune!90!black},
  milestone/.append style={fill=red!70,draw=red!80!black},
  title/.append style={fill=anbleu!12},
  title label font=\footnotesize\bfseries,
  bar label font=\footnotesize,
  group label font=\footnotesize\bfseries,
  milestone label font=\scriptsize,
  y unit chart=0.55cm, y unit title=0.6cm
  ]{1}{12}
  \gantttitle{Mois 1}{4}\gantttitle{Mois 2}{4}\gantttitle{Mois 3}{4}\\
  \gantttitlelist{1,...,12}{1}\\

  \ganttgroup{Cadrage \& Analyse}{1}{5}\\
  \ganttbar{Planification \& cahier des charges}{1}{2}\\
  \ganttbar{Diag. contexte / package / classes}{2}{4}\\
  \ganttbar{Cas d'utilisation \& activités}{3}{5}\\
  \ganttbar{Diag. séquence système}{4}{5}\\
  \ganttmilestone{Cahier d'analyse validé}{5}\\

  \ganttgroup{Conception}{5}{8}\\
  \ganttbar{Séquences \& classes techniques}{5}{7}\\
  \ganttbar{Modèle de données \& architecture}{6}{8}\\
  \ganttmilestone{Conception figée}{8}\\

  \ganttgroup{Développement}{6}{11}\\
  \ganttbar{Backend Spring Boot (API, sécurité JWT)}{6}{11}\\
  \ganttbar{Frontend Next.js (pages, intégration)}{7}{11}\\

  \ganttgroup{Qualité \& Livraison}{10}{12}\\
  \ganttbar{Tests (unitaires, intégration, E2E)}{10}{12}\\
  \ganttbar{Recette (UAT) \& corrections}{11}{12}\\
  \ganttbar{Déploiement \& documentation}{12}{12}\\
  \ganttmilestone{Livraison finale}{12}
\end{ganttchart}}
\caption{Diagramme de Gantt du projet \textsc{Assurance Nation} : de la conception
(diagrammes UML) au code et au déploiement, sur 3 mois.}
\end{figure}

%==============================================================================
%  SECTION 11
%==============================================================================
\section{Cadre Contractuel, Financier et Juridique}

\emph{Le présent cadre est établi conformément au droit camerounais applicable.}

\subsection{Structure de Tarification et Propriété Intellectuelle}
\textbf{Modèle économique :} dans le contexte académique, la prestation est réalisée
au \textbf{forfait} (engagement de résultat sur le périmètre du §2). Pour une
exploitation réelle, trois modèles sont envisageables : développement au forfait
(prix fixe), facturation en régie (au temps passé), ou abonnement récurrent (SaaS).

\textbf{Propriété intellectuelle :} conformément à la \textbf{loi camerounaise
n\textsuperscript{o}~2000/011 du 19 décembre 2000} relative au droit d'auteur et aux
droits voisins, et aux dispositions de l'\textbf{OAPI} (Accord de Bangui révisé) :
\begin{itemize}[leftmargin=1.4em]
  \item le code source développé spécifiquement fait l'objet d'une \textbf{cession
  totale des droits patrimoniaux} au profit du commanditaire (l'organisme de sécurité
  sociale), une fois la prestation soldée ;
  \item les modules et bibliothèques tiers (open source) demeurent sous leurs licences
  respectives, en \textbf{concession d'usage} ;
  \item des \textbf{clauses de confidentialité} protègent le patrimoine informationnel
  (données médicales, savoir-faire) partagé durant l'ingénierie.
\end{itemize}

\subsection{Pénalités, Responsabilités et Résolution des Litiges}
\begin{itemize}[leftmargin=1.4em]
  \item \textbf{Pénalités de retard} : en cas de dépassement abusif des échéances
  contractuelles, application d'une pénalité forfaitaire par jour ouvré de retard,
  plafonnée, telle que stipulée au contrat.
  \item \textbf{Non-livraison / non-conformité} : la non-livraison de fonctionnalités
  classées \emph{Must} ou le non-respect du SLA de disponibilité (cf. §5) ouvrent droit
  à pénalités et, le cas échéant, à la résiliation.
  \item \textbf{Responsabilité} : les obligations et la limitation de responsabilité
  s'apprécient au regard du \textbf{Code civil} applicable au Cameroun ; le prestataire
  est tenu d'une obligation de moyens renforcée sur la sécurité des données.
  \item \textbf{Escalade et arbitrage} : tout différend est d'abord soumis à un
  règlement amiable ; à défaut, il relève des juridictions camerounaises compétentes,
  l'arbitrage pouvant être conduit selon le \textbf{droit OHADA} (Acte uniforme relatif
  à l'arbitrage, CCJA).
\end{itemize}

%==============================================================================
%  SECTION 12
%==============================================================================
\section{Annexes et Matrice de Traçabilité (RTM)}

\subsection{Annexes — Artefacts Visuels}
Sont annexés au présent cahier des charges, en complément du \emph{Cahier d'Analyse} :
\begin{itemize}[leftmargin=1.4em]
  \item les \textbf{diagrammes UML} (contexte, packages, classes métier et techniques,
  objets, cas d'utilisation, activités, séquences système et techniques) ;
  \item le \textbf{schéma d'architecture de déploiement} (cf. §7.2) ;
  \item le \textbf{schéma du modèle de données} relationnel (cf. §4.8) ;
  \item les \textbf{wireframes} et maquettes d'interface des écrans principaux
  (connexion, tableau de bord, consultations, remboursements).
\end{itemize}

\subsection{Matrice de Traçabilité des Exigences (RTM)}
La matrice ci-dessous relie chaque \textbf{besoin métier} initial à son
\textbf{exigence logicielle} détaillée et au \textbf{scénario de test} qui la valide,
permettant des analyses d'impact fiables en cas de demande de changement.

\begin{center}
\begin{longtable}{|L{3.6cm}|L{2.4cm}|L{2.2cm}|L{4.6cm}|}
\hline
\rowcolor{anbleu!12}\textbf{Besoin métier} & \textbf{Exigence} & \textbf{Priorité} & \textbf{Scénario de test}\\\hline
\endhead
Centraliser l'accès sécurisé des acteurs & EF-AUTH-01/02/03 & Must & TST-AUTH : inscription, connexion, refresh, déconnexion.\\\hline
Gérer les assurés et leurs médecins traitants & EF-USR-01/03 & Must/Should & TST-USR-ASSURE : CRUD, unicité NSS, affectation médecin traitant.\\\hline
Gérer le référentiel des médecins & EF-USR-02 & Must & TST-USR-MED : CRUD, unicité RPPS, spécialité.\\\hline
Permettre la gestion de compte personnel & EF-USR-04 & Must & TST-PROFIL : mise à jour profil, changement de mot de passe.\\\hline
Tracer le parcours de soins (consultation) & EF-CONS-01/02 & Must & TST-CONS : création + feuille de maladie, visibilité par rôle, annulation.\\\hline
Prescrire médicaments et orientations & EF-PRESC-01/02 & Must & TST-PRESC : médicament (durée), consultation spécialiste (motif).\\\hline
Consulter les feuilles de maladie & EF-FM-01 & Must & TST-FM : filtres période/statut, lecture seule si remboursée.\\\hline
Automatiser et fiabiliser les remboursements & EF-RB-01/02/03 & Must & TST-RB : calcul 100/80~\%, workflow d'états, justificatif PDF.\\\hline
Piloter l'activité (décision) & EF-RB-04 & Should & TST-DASH : compteurs, montants, graphiques.\\\hline
Garantir la conformité et la traçabilité & EF-AUD-01 & Must & TST-AUDIT : journalisation des opérations sensibles.\\\hline
Informer les assurés & EF-NOT-01 & Should & TST-NOTIF : envoi de courriels, non-blocage en cas d'échec.\\\hline
\end{longtable}
\end{center}

\end{document}
